\section{Three harnesses and one instruction file}
\label{sec:systems}

% Remaster, 3 October 2026.

The three harnesses are three design patterns of agentic AutoML: a workflow
with fixed stages that the model configures, a team of specialised agents,
and a single agent in an open execution loop~\citep{anthropic2024agents}.
Figure~\ref{fig:systems} draws them and Table~\ref{tab:designpoints} lists
what each one fixes in advance and what it leaves to the model.

\subsection{FEDOT.LLM: a fixed-stage workflow}

FEDOT.LLM~\citep{fedotllm} pairs an LLM agent with the FEDOT
framework~\citep{nikitin2022fedot} in four fixed stages. Data understanding
extracts paths, column types, target, metric and submission format;
configuration and code generation produce a FEDOT configuration and fill a
code template with fixed input and output contracts; an execution loop runs
the code in a sandbox and patches it until it succeeds or a retry cap is
reached; reporting saves the script, the trained pipeline and the
predictions. Two properties matter for the argument. The modelling tool is
fixed by design: every task is solved with FEDOT, which searches over model
pipelines and tunes them within a time budget. And the workflow accepts
tabular inputs only.

\subsection{AutoDS-Tools: a multi-agent system}

AutoDS-Tools~\citep{autodstools} coordinates six agents through written
reports. An Analyst explores the data, a Researcher studies the target
library and writes usage notes, a Planner turns both into an execution plan,
a Coder writes and refines the solution against execution results, a
Debugger absorbs execution errors outside the main context, and a Presenter
audits the result. Each stage writes a Markdown report that the next one
reads. The system may use any library installed in its container and may
restructure the solution between iterations.

\subsection{Terminus-2: a minimal single-agent harness}

Terminus-2 is the reference agent of Terminal-Bench~\citep{merrill2026terminalbench}.
It gives the model one tool, a shell in a working directory, and an
execution loop: read the task statement, inspect the data, write code, run
it, read the output, repeat until a submission is written or the budget is
exhausted. It prescribes nothing about how the task should be solved.

\begin{table}[t]
\centering
\caption{What each harness fixes in advance and what it leaves to the model.
The last two rows say where the instruction file of Section~\ref{sec:file}
can be switched.}
\label{tab:designpoints}
\footnotesize
\setlength{\tabcolsep}{3.5pt}
\begin{tabular}{llll}
\toprule
 & FEDOT.LLM & AutoDS-Tools & Terminus-2 \\
\midrule
Design pattern     & workflow        & multi-agent     & single agent \\
Modelling tool     & FEDOT, fixed     & free            & free \\
Plan               & fixed stages    & generated       & emergent \\
Execution loop     & bounded retries & per-agent       & open \\
Input modalities   & tabular only    & any             & any \\
Left to the model  & configuration   & everything      & everything \\
\midrule
Library instruction    & built in        & on / off        & off, appendable \\
Discipline instruction & absent          & on / off        & off, appendable \\
\bottomrule
\end{tabular}
\end{table}

% Figure 1: the three systems on one schematic.  Redrawn 21 September 2026
% from the prototype in review_protos/fig1_proto.tex.  Colours cFedot, cAuto,
% cTerm, cLayer, ink, muted, ctx are defined in main.tex and match
% make_figures.py.  Scores and ranks come from tables/numbers.tex.
\begin{figure*}[t]
\centering
\resizebox{\textwidth}{!}{%
\begin{tikzpicture}[x=1mm,y=1mm, font=\footnotesize, text=ink,
  stage/.style = {draw=ctx, fill=black!10, rounded corners=1.2pt, line width=0.4pt,
                  minimum height=9mm, minimum width=25mm, text width=23.5mm,
                  align=center, inner xsep=0.5pt, inner ysep=1.5pt},
  free/.style  = {stage, fill=white, line width=0.9pt},
  chip/.style  = {fill=cLayer!60, rounded corners=1pt, inner xsep=2.2pt, inner ysep=1.6pt,
                  font=\scriptsize, align=center},
  io/.style    = {inner sep=1pt, align=center},
  arr/.style   = {-{Latex[length=1.5mm,width=1.2mm]}, draw=ink!70, line width=0.5pt},
  loop/.style  = {arr, draw=muted, rounded corners=1pt},
  env/.style   = {rounded corners=2pt, draw=ctx, densely dotted, line width=0.45pt},
  lbl/.style   = {font=\scriptsize, text=muted, inner sep=1pt, align=center},
  lane/.style  = {anchor=east, align=right, inner sep=0pt},
]
% ---- geometry ----------------------------------------------------------
\def\xLane{21}     % right edge of lane labels
\def\xTask{28}     % centre of the input node
\def\xA{34}\def\xB{61}\def\xC{88}\def\xD{115}   % left edges of the 4 stage columns
\def\wS{25}        % stage width
\def\xOut{149}     % centre of the output node
\def\xBar{160}     % score bar origin
\def\wBar{18}      % score bar span for score = 1
\def\yA{0}\def\yB{-26}\def\yC{-52}
\def\hRow{9}

% ---- header ------------------------------------------------------------
\node[lbl, anchor=south] at (\xTask,13) {input};
\node[lbl, anchor=south] at ({(\xA+\xD+\wS)/2},13) {stages};
\node[lbl, anchor=south] at (\xOut,13) {output};
\node[lbl, anchor=south, text width=30mm] at ({\xBar+\wBar/2},13)
     {TabReD score\\(0 weakest, 1 strongest)};
\draw[ctx, line width=0.4pt] (\xBar,11.5) -- ++(\wBar,0);
\foreach \t/\l in {0/0, 1/1} \draw[ctx, line width=0.4pt] ({\xBar+\t*\wBar},11.5) -- ++(0,-1)
   node[lbl, below, inner sep=0.3pt] {\l};

% ---- helpers -----------------------------------------------------------
% loop above the span [#1,#2] in row #3, label #4
\newcommand{\iterloop}[4]{%
  \draw[loop] ({#2-3},{#3+\hRow/2}) -- ({#2-3},{#3+\hRow/2+2.6}) -- ({#1+3},{#3+\hRow/2+2.6}) -- ({#1+3},{#3+\hRow/2+0.3});
  \node[lbl, anchor=south] at ({(#1+#2)/2},{#3+\hRow/2+2.9}) {#4};}
% dotted container frame around the span [#1,#2] in row #3
\newcommand{\envframe}[3]{%
  \begin{scope}[on background layer]
    \draw[env] ({#1-1.4},{#3-\hRow/2-3.4}) rectangle ({#2+1.4},{#3+\hRow/2+6.4});
  \end{scope}
  \node[lbl, anchor=south east, inner sep=1.2pt] at ({#2+1.4},{#3-\hRow/2-3.4}) {container};}
% score bar in row #1, value #2, colour #3, rank #4
\newcommand{\scorebar}[4]{%
  \fill[#3] (\xBar,{#1-1.4}) rectangle ({\xBar+#2*\wBar},{#1+1.4});
  \node[anchor=west, inner sep=1.5pt, align=left] at ({\xBar+#2*\wBar},#1)
       {#2\\[-1pt]{\scriptsize\color{muted}rank #4 of \nMethods}};}

% ============================ row 1: FEDOT.LLM ==========================
\node[lane, text=cFedot] at (\xLane,\yA) {FEDOT.LLM\\[-1pt]{\scriptsize\color{muted}all 4 stages fixed}};
\node[io] (t1) at (\xTask,\yA) {task\\[-2pt]{\scriptsize\color{muted}tabular}};
\node[stage, anchor=west] (a1) at (\xA,\yA) {data\\understanding};
\node[stage, anchor=west] (b1) at (\xB,\yA) {FEDOT config\\\phantom{x}};
\node[chip, anchor=south] at ([yshift=0.9mm]b1.south) {\Ktool{} welded in};
\node[stage, anchor=west] (c1) at (\xC,\yA) {code from\\template};
\node[stage, anchor=west] (d1) at (\xD,\yA) {run $+$\\bounded repair};
\node[io] (e1) at (\xOut,\yA) {submission};
\foreach \i/\j in {t1/a1,a1/b1,b1/c1,c1/d1,d1/e1} \draw[arr] (\i) -- (\j);
\iterloop{\xD}{\xD+\wS}{\yA}{until pass or retry cap}
\envframe{\xC}{\xD+\wS}{\yA}
\scorebar{\yA}{\nFedot}{cFedot}{\rkFedot}

% ============================ row 2: AutoDS-Tools ========================
\node[lane, text=cAuto] at (\xLane,\yB) {AutoDS-Tools\\[-1pt]{\scriptsize\color{muted}2 of 4 stages fixed}};
\node[io] (t2) at (\xTask,\yB) {task\\[-2pt]{\scriptsize\color{muted}any}};
\node[stage, anchor=west] (a2) at (\xA,\yB) {analyst,\\researcher};
\node[free, draw=cAuto, anchor=west] (b2) at (\xB,\yB) {planner:\\plan $+$ spec};
\node[free, draw=cAuto, anchor=west] (c2) at (\xC,\yB) {coder\\any library};
\node[stage, anchor=west] (d2) at (\xD,\yB) {debugger,\\presenter};
\node[io] (e2) at (\xOut,\yB) {submission};
\foreach \i/\j in {t2/a2,a2/b2,b2/c2,c2/d2,d2/e2} \draw[arr] (\i) -- (\j);
\iterloop{\xC}{\xD+\wS}{\yB}{per-agent retries}
\envframe{\xC}{\xD+\wS}{\yB}
\node[chip, anchor=north] (k2) at (\xTask,{\yB-6.5}) {\Ktool{} $+$ \Kdisc\\[-1pt]on / off};
\draw[arr, draw=muted] (k2.north) -- (t2.south);
\node[lbl, anchor=west, align=left] at ({\xA+2},{\yB-9.1}) {composed on the host into the instruction that every agent reads};
\scorebar{\yB}{\nAutoDS}{cAuto}{\rkAutoDS}
\draw[white, line width=0.6pt] ({\xBar+\nAutoDSOff*\wBar},{\yB-1.4}) -- ++(0,2.8);
\node[lbl, anchor=south] at ({\xBar+\nAutoDSOff*\wBar},{\yB+1.8}) {\nAutoDSOff{} with the layer off};

% ============================ row 3: Terminus-2 ==========================
\node[lane, text=cTerm] at (\xLane,\yC) {Terminus-2\\[-1pt]{\scriptsize\color{muted}no fixed stage}};
\node[io] (t3) at (\xTask,\yC) {task file\\[-2pt]{\scriptsize\color{muted}any}};
\node[free, draw=cTerm, anchor=west, minimum width={(\xD+\wS-\xA)*1mm}, text width={(\xD+\wS-\xA-3)*1mm}] (a3) at (\xA,\yC)
     {shell $+$ working directory $+$ execution loop\\read data, write code, run, inspect, repeat: every modelling decision is taken here};
\node[io] (e3) at (\xOut,\yC) {submission};
\draw[arr] (t3) -- (a3); \draw[arr] (a3) -- (e3);
\iterloop{\xA}{\xD+\wS}{\yC}{until submission or budget}
\envframe{\xA}{\xD+\wS}{\yC}
\node[chip, anchor=north] (k3) at (\xTask,{\yC-6.5}) {\Ktool{} $+$ \Kdisc\\[-1pt]on / off};
\draw[arr, draw=muted] (k3.north) -- (t3.south);
\node[lbl, anchor=west, align=left] at ({\xA+2},{\yC-9.1}) {appended to the task file};
\scorebar{\yC}{\nTerminus}{cTerm}{\rkTerminus}

% ---- legend --------------------------------------------------------------
\begin{scope}[shift={(\xA,-67)}]
  \node[stage, minimum width=6mm, minimum height=3.6mm, text width=5mm] at (0,0) {};
  \node[anchor=west] at (3.5,0) {role fixed by the architecture};
  \node[free, draw=ink!70, minimum width=6mm, minimum height=3.6mm, text width=5mm] at (54,0) {};
  \node[anchor=west] at (57.5,0) {the model takes the modelling decisions here};
  \node[chip, minimum width=6mm, minimum height=3.6mm] at (126,0) {};
  \node[anchor=west] at (129.5,0) {prescribed-knowledge layer};
\end{scope}
\end{tikzpicture}}
\caption{The three systems on one schematic. Grey stages have their role fixed
by the architecture; outlined stages are where the model takes the modelling
decisions. The yellow chip marks where the prescribed-knowledge layer enters:
welded into the FEDOT configuration, composed into the instruction that every
AutoDS-Tools agent reads, appended to the Terminus-2 task file. The loop above
a row shows what repeats and what ends it; the dotted frame is the container.
The right column gives the mean normalised TabReD score of
Section~\ref{sec:ceiling}, where 0 is the weakest and 1 the strongest of the
eighteen tuned methods on each task, and the average rank among the
\nMethods{} methods of Table~\ref{tab:leaderboard}. Fixed structure falls from
top to bottom; the score peaks in the middle row.}
\label{fig:systems}
\end{figure*}


\subsection{The instruction file and its two parts}
\label{sec:file}

AutoDS-Tools ships with an expert instruction file that is composed on the
host and prepended to the task statement every agent reads. In current
terminology it is an agent skill~\citep{li2026skillsbench}: standing
procedural instructions, independent of the task. It has two parts, which
we measure separately.
\begin{description}
  \item[Library instruction.] Which library to use for the task family: a
        dedicated tabular AutoML library (LightAutoML~\citep{vakhrushev2021lightautoml})
        for tabular data, a vision, graph or transformer library for the other
        modalities. It states that the library is pre-installed and forbids a
        fallback to general-purpose estimators.
  \item[Discipline instruction.] How to train and how to check oneself:
        the validation protocol, handling of temporal shift, leakage checks,
        comparison against a trivial predictor, allocation of the time budget,
        the prohibition of degenerate submissions, and advice on training
        time.
\end{description}
On a tabular task the file is about 5$\times$ as long as the task statement,
and the library instruction is about a quarter of it. An ablation that
switches the whole file at once cannot tell the two parts apart.

In FEDOT.LLM the library instruction is the workflow itself: every task goes
through FEDOT. Terminus-2 ships no instruction of its own, so we test the
file in a second harness by appending it to the task statement
(Table~\ref{tab:designpoints}).
